\section{Данные}
\label{sec:data}

Пять таблиц, около 35 тысяч строк. Устройство этих данных важнее их объёма: вся конструкция модели вырастает из одной асимметрии.
Обучающая матрица разрежена и смещена отбором, тестовая плотная.

\begin{table}[htbp]\centering\small
\begin{tabularx}{\textwidth}{@{}Xllp{62mm}@{}}
\toprule
Источник & Форма & Заполнено & Что это \\
\midrule
Кривые доза-эффект, прямое плечо & $4905\times 4$ & 26--48\,\% & $\pic$ плюс границы интервала и стандартная ошибка \\
Кривые, плечо с преинкубацией & $6145\times 4$ & 21--58\,\% & то же для условия TDI \\
Emax, оба плеча & $6146\times 8$ & частично & предельная глубина подавления и её интервал \\
Одноточечный скрининг & $4376\times 4$ & 100\,\% & $\log_2$fc при 49.5 мкМ и стандартная ошибка \\
Тест & $750\times 4$ & --- & слепой, плотный \\
\bottomrule
\end{tabularx}
\caption{Пять источников. Пересечения между ферментами внутри таблицы кривых маленькие: одновременно
для 1A2 и 2C9 измерено 295 соединений, для 2C9 и 2D6 --- 230. Скрининг --- единственный плотный
источник, и он покрывает весь диапазон активностей, включая неактивные соединения, которых в кривых
почти нет.}
\end{table}

\subsection{Что меряют эти таблицы}

Кривая доза-эффект: соединение разводят по концентрациям, для каждой измеряют, насколько подавлена
активность фермента, и через полученные точки проводят сигмоиду. Из подгонки получают $\pic$ ---
отрицательный десятичный логарифм концентрации, при которой активность падает вдвое. Байесовская подгонка
даёт не только точку, но и интервал, отражающий, насколько уверенно эта точка определена. У соединений,
чья кривая не дошла до плато, интервал широкий.

Одноточечный скрининг: то же измерение, но при одной концентрации 49.5 мкМ. На выходе $\log_2$fc ---
логарифм отношения сигнала с соединением к сигналу без него. Ноль означает, что ничего не произошло,
отрицательные значения --- что активность подавлена.

Полные кривые снимали не для всех подряд. Соединение поднимали на подробное
измерение, если оно показало эффект на скрининге. Это отбор по зависимой переменной, и он деформирует
обучающую выборку способом, который разобран в разделе~\ref{sec:d6}.

Правило это не универсальное. Скрининг есть у 4375 из 4905 соединений с кривыми; у оставшихся 530 его
нет, и все 530 несут метку \emph{только} по 3A4 --- на 1A2, 2C9 и 2D6 у них меток ноль. То есть для трёх
ферментов отбор по скринингу описывает выборку целиком, а на 3A4 почти четверть кривых (530 из 2335)
пришла каким-то другим путём, и для них аргумент про отбор просто не про них.

\subsection{Как собран тестовый набор}

Организаторы описали процедуру: взяли по 25 лучших ингибиторов каждого из 1A2, 2C9 и 3A4, к каждому
добавили по 10 ближайших соседей по фингерпринту, итого $3\times 25\times 10 = 750$.

Это проверяется, и проверка сходится частично. Здесь важно, что порог кластеризации выбирается, а не
выводится, и картина от него зависит сильно. При рабочем пороге 0.35 из 750 тестовых соединений выходит
477 групп с медианным размером один, групп по пять и более членов 21, и в них лежат 155 соединений.
При 0.48 --- 203 группы с медианным размером два, 56 групп по пять и более, в них 496 соединений.
То есть серийная структура в наборе есть, но её видимый охват меняется от пятой части до двух третей
в зависимости от того, что считать серией. Ровно 75 серий ровно по 10 не восстанавливается ни при каком
пороге --- число 75 получается делением 750 на 10.

Что восстанавливается уверенно, так это положение якорей: оно от порога почти не зависит. Для каждой
группы найдено ближайшее обучающее соединение, и его активность стоит на 92--94 перцентиле обучающего
распределения по 1A2 и на 97.5--99 по 2C9 при любом пороге от 0.35 до 0.50, тогда как по 2D6, который
в отборе не участвовал, --- на 58--66. По 3A4 оценка гуляет заметно шире, от 82 до 95, и опирать на неё
точное число нельзя. Оговорка, без которой эти проценты читаются неверно: якорей с меткой всего от
девяти до тридцати семи на фермент, так что медиана здесь имеет широкий интервал. Медианное сходство
Танимото тестового соединения с ближайшим обучающим равно 0.587.

Значит распределение меток на тесте не совпадает с обучающим: тест обогащён активными по трём ферментам
и близок к фону по четвёртому. И устроен он как набор серий аналогов вокруг известных хитов, а не как
случайная выборка из химического пространства. Проверять модель надо соответственно.
